% Generated draft: requires mathematical and source review.
\documentclass[11pt]{article}
\usepackage[margin=25mm]{geometry}
\usepackage{fontspec}
\setmainfont{Latin Modern Roman}
\usepackage{amsmath,amssymb,mathtools,unicode-math}
\setmathfont{Latin Modern Math}
\usepackage{enumitem,needspace,xurl,hyperref}
\setlength{\parindent}{0pt}
\setlength{\parskip}{5pt}
\setlength{\emergencystretch}{3em}
\newcommand{\N}{\mathbb{N}}
\newcommand{\Z}{\mathbb{Z}}
\newcommand{\Q}{\mathbb{Q}}
\newcommand{\R}{\mathbb{R}}
\newcommand{\C}{\mathbb{C}}
\newcommand{\F}{\mathbb{F}}
\newcommand{\A}{\mathbb{A}}
\newcommand{\PP}{\mathbb{P}}
\newcommand{\E}{\mathbb{E}}
\newcommand{\eps}{\varepsilon}
\newcommand{\dd}{\,\mathrm{d}}
\begin{document}
\section*{UnsolvedMath ID: 30003353}
\textbf{Source problem code:} catalogue entry
\subsection{Definitions and mathematical statement}
Let \(n\in\Z\) with \(n\geq 2\). A smooth homotopy sphere of dimension \(4n-1\) is a closed smooth manifold \(\Sigma\) homotopy equivalent to \(S^{4n-1}\). Define
\[
\mathcal R_{\sec\geq 0}(\Sigma)
=\{g\mid g\text{ is a smooth Riemannian metric on }\Sigma,\quad
\sec_g(\sigma)\geq 0\text{ for every }p\in\Sigma
\text{ and every two-dimensional subspace }\sigma\subset T_p\Sigma\}.
\]
The hypothesis that \(\Sigma\) is nonnegatively curved means that \(\mathcal R_{\sec\geq 0}(\Sigma)\neq\varnothing\).

For a precise conventional interpretation, give this metric space the \(C^\infty\) topology and define
\[
\mathcal M_{\sec\geq 0}(\Sigma)
=\mathcal R_{\sec\geq 0}(\Sigma)/\operatorname{Diff}(\Sigma),
\]
where the full group of smooth diffeomorphisms acts by pullback and the orbit space has the quotient topology. The supplied record does not explicitly specify this topology or the diffeomorphism group; these conventions require confirmation against the source.

\textbf{Question.} For every integer \(n\geq 2\) and every smooth homotopy sphere \(\Sigma\) of dimension \(4n-1\) admitting a metric of nonnegative sectional curvature, does \(\mathcal M_{\sec\geq 0}(\Sigma)\) have infinitely many connected components?

Here components mean connected components, as in the supplied cleaned statement; no identification with path components is assumed.
\subsection{Short English statement}
Does every smooth homotopy sphere of dimension \(4n-1\geq 7\) that admits nonnegative sectional curvature have infinitely many connected components in its moduli space of such metrics? The formulation uses the \(C^\infty\) topology and quotient by all diffeomorphisms, subject to source confirmation.
\subsection{Sources}
\begin{itemize}
\item Supplied local UnsolvedMath registry record, permanent ID \(30003353\), source problem code \texttt{catalogue entry}, titled \emph{Components of Nonnegative Curvature Metrics on Homotopy Spheres}, in the Oberwolfach Reports Open Problems collection. No individual proposer is supplied; the proposed year is recorded as 2017.
\item Supplied report citation: \emph{Mini-Workshop: Spaces and Moduli Spaces of Riemannian Metrics (2018)}, Oberwolfach Reports, pp.\ 133--166, DOI 10.4171/owr/2017/3. \url{https://doi.org/10.4171/owr/2017/3}. The background describes this report as documenting workshop discussions and open problems. The supplied title year and recorded proposed year differ; the bibliographic details require verification.
\item Supplied literature source: \emph{Sphere bundles and nonnegative sectional moduli}. \url{https://doi.org/10.1007/s10711-022-00713-7}.
\item Supplied literature source: \emph{Moduli spaces of nonnegatively curved metrics on exotic spheres}. \url{https://arxiv.org/abs/2107.02310}.
\item Supplied literature assessment, checked on \textbf{2026-08-21}: the record labels the problem \emph{partially solved} and reports affirmative answers for broad classes of homotopy spheres, including seven-dimensional Milnor spheres and later large classes of manifolds homeomorphic to the seven-sphere, while reporting that the full dimension-general statement remains open. This is supplied dated curation derived from the OWR material, not an independently verified literature assessment; its status and cited claims require independent verification.
\end{itemize}
\end{document}
